\makeatletter

\newtcolorbox{missionbox}[1][]{
  enhanced,
  breakable,
  colback=cqpanel,
  colframe=cqcyan!80!black,
  coltitle=cqink,
  coltext=cqink,
  fonttitle=\bfseries,
  fontupper=\small,
  title={#1},
  boxrule=0.9pt,
  arc=2mm,
  left=2.5mm,
  right=2.5mm
}

\newtcolorbox{conceptbox}[1][]{
  enhanced,
  breakable,
  colback=white,
  colframe=cqlime!60!black,
  coltitle=cqink,
  coltext=cqink,
  fonttitle=\bfseries,
  fontupper=\small,
  title={#1},
  boxrule=0.9pt,
  arc=2mm
}

\newtcolorbox{examplebox}[1][]{
  enhanced,
  breakable,
  colback=cqpanel,
  colframe=cqyellow!55!black,
  coltitle=cqink,
  coltext=cqink,
  fonttitle=\bfseries,
  fontupper=\small,
  title={#1},
  boxrule=0.9pt,
  arc=2mm
}

\newtcolorbox{guidebox}[1][]{
  enhanced,
  breakable,
  colback=white,
  colframe=cqyellow!70!black,
  coltitle=cqink,
  coltext=cqink,
  fonttitle=\bfseries,
  fontupper=\small,
  title={#1},
  boxrule=0.9pt,
  arc=2mm
}

\newtcolorbox{semibox}[1][]{
  enhanced,
  breakable,
  colback=white,
  colframe=cqcyan!80!black,
  coltitle=cqink,
  coltext=cqink,
  fonttitle=\bfseries,
  fontupper=\small,
  title={#1},
  boxrule=0.9pt,
  arc=2mm
}

\newtcolorbox{solobox}[1][]{
  enhanced,
  breakable,
  colback=cqpanel,
  colframe=cqlime!55!black,
  coltitle=cqink,
  coltext=cqink,
  fonttitle=\bfseries,
  fontupper=\small,
  title={#1},
  boxrule=0.9pt,
  arc=2mm
}

\newtcolorbox{machinebox}[1][]{
  enhanced,
  breakable,
  colback=cqcodebg,
  colframe=cqcyan,
  coltitle=white,
  fonttitle=\bfseries,
  fontupper=\small\color{white},
  title={#1},
  boxrule=0.9pt,
  arc=2mm
}

\newtcolorbox{checkpointbox}[1][]{
  enhanced,
  breakable,
  colback=cqpanel,
  colframe=cqline,
  coltitle=cqink,
  coltext=cqink,
  fonttitle=\bfseries,
  fontupper=\small,
  title={#1},
  boxrule=0.9pt,
  arc=2mm
}

\newtcolorbox{languagebox}[1][]{
  enhanced,
  breakable,
  colback=white,
  colframe=cqdarkpanel,
  coltitle=cqink,
  coltext=cqink,
  fonttitle=\bfseries,
  fontupper=\small,
  title={#1},
  boxrule=0.9pt,
  arc=2mm
}

\newtcolorbox{bossbox}[1][]{
  enhanced,
  breakable,
  colback=cqcodebg,
  colframe=cqyellow,
  coltitle=white,
  fonttitle=\bfseries,
  fontupper=\small\color{white},
  title={#1},
  boxrule=1pt,
  arc=2mm
}

\newcommand{\IssueBanner}[2]{%
  \begin{tcolorbox}[
    enhanced,
    breakable,
    colback=cqpanel,
    colframe=cqcyan!75!black,
    boxrule=0.9pt,
    arc=2mm,
    left=2.5mm,
    right=2.5mm,
    top=1.5mm,
    bottom=1.6mm,
    overlay={
      \fill[cqtoolbar] ([yshift=-6.8mm]frame.north west) rectangle (frame.north east);
      \node[anchor=west, text=cqcyan, font=\ttfamily\scriptsize\bfseries]
        at ([xshift=2mm,yshift=-3.4mm]frame.north west) {code quest | logic map};
      \node[anchor=east, text=cqmuted, font=\ttfamily\scriptsize]
        at ([xshift=-2mm,yshift=-3.4mm]frame.north east) {study mode};
    }
  ]
    \vspace{2.7mm}
    {\ttfamily\footnotesize\color{cqmagenta}>\_ urmarim logica jocului}\par
    \vspace{0.15em}
    {\LARGE\bfseries\color{cqink} #1}\par
    \vspace{0.1em}
    {\normalsize\color{cqdarkpanel} #2}
  \end{tcolorbox}
}

\newcommand{\PageHeader}[3]{%
  \begin{tcolorbox}[
    enhanced,
    breakable,
    colback=cqpanel,
    colframe=cqcyan!70!black,
    boxrule=0.85pt,
    arc=2mm,
    left=2.5mm,
    right=2.5mm,
    top=1.4mm,
    bottom=1.4mm,
    overlay={
      \fill[cqtoolbar] ([yshift=-6.3mm]frame.north west) rectangle (frame.north east);
      \node[anchor=west, text=cqcyan, font=\ttfamily\scriptsize\bfseries]
        at ([xshift=2mm,yshift=-3.1mm]frame.north west) {code quest};
    }
  ]
    \vspace{2.4mm}
    {\Large\bfseries\color{cqink} #1}\par
    \vspace{0.12em}
    {\small\color{cqdarkpanel} #2}
  \end{tcolorbox}
}

\newcommand{\FieldLine}[1]{%
  \par\noindent\textbf{#1}\dotfill\par
}

\newcount\cqanswercount
\newcommand{\AnswerLine}{\par\noindent\color{cqline}\rule{\linewidth}{0.5pt}\par\vspace{0.7em}}
\newcommand{\AnswerSpace}[1]{%
  \cqanswercount=0
  \loop\ifnum\cqanswercount<#1
    \advance\cqanswercount by 1
    \AnswerLine
  \repeat
}

\newcommand{\QRCodePlaceholder}{%
  \begin{tikzpicture}[x=1.7mm,y=1.7mm]
    \draw[cqcyan, line width=0.6pt] (0,0) rectangle (14,14);
    \fill[cqcodebg] (1,1) rectangle (5,5);
    \fill[white] (1.8,1.8) rectangle (4.2,4.2);
    \fill[cqcodebg] (2.5,2.5) rectangle (3.5,3.5);
    \fill[cqcodebg] (9,1) rectangle (13,5);
    \fill[white] (9.8,1.8) rectangle (12.2,4.2);
    \fill[cqcodebg] (10.5,2.5) rectangle (11.5,3.5);
    \fill[cqcodebg] (1,9) rectangle (5,13);
    \fill[white] (1.8,9.8) rectangle (4.2,12.2);
    \fill[cqcodebg] (2.5,10.5) rectangle (3.5,11.5);
    \fill[cqcodebg] (7,7) rectangle (10,10);
    \fill[cqcodebg] (6,4) rectangle (7,5);
    \fill[cqcodebg] (12,9) rectangle (13,10);
  \end{tikzpicture}
}

\newcommand{\BlocklyExperiment}[3]{%
  \begin{tcolorbox}[
    enhanced,
    colback=white,
    colframe=cqmagenta,
    boxrule=1pt,
    arc=2mm,
    title={Experiment Blockly: #2}
  ]
    \small #3
    \begin{flushright}
      \begin{tikzpicture}
        \node[anchor=center] at (0,0) {\QRCodePlaceholder};
        \node[anchor=north, font=\ttfamily\tiny, cqmuted] at (0,-1.3) {#1};
      \end{tikzpicture}
    \end{flushright}
  \end{tcolorbox}
}

\newcommand{\SimulatorCard}[4]{%
  \begin{tcolorbox}[
    enhanced,
    breakable,
    colback=white,
    colframe=cqcyan,
    title={Simulator interactiv: #2}
  ]
  \small
  \textbf{Ce arata simulatorul:} #3

  \textbf{Muta controlul:} sliderul principal al simulatorului si observa modificarile in timp real.

  \textbf{Inainte sa scanezi (predictie):} Ce crezi ca se va intampla daca modifici valorile extreme?

  \textbf{Incearca si observa:} #4

  \begin{flushright}
    \begin{tikzpicture}
      \node[anchor=center] at (0,0) {\QRCodePlaceholder};
      \node[anchor=north, font=\ttfamily\tiny, cqmuted] at (0,-1.3) {#1};
    \end{tikzpicture}
  \end{flushright}
  \end{tcolorbox}
}

\expandafter\newcommand\csname cqSyntaxRows@variabile\endcsname{%
  scor <- 0 & int scor = 0; & scor = 0 & int scor = 0; \\
  vieti <- 3 & int vieti = 3; & vieti = 3 & int vieti = 3; \\%
}

\expandafter\newcommand\csname cqSyntaxRows@operatori\endcsname{%
  scor <- scor + 1 & scor = scor + 1; & scor = scor + 1 & scor = scor + 1; \\
  vieti <- vieti - 1 & vieti = vieti - 1; & vieti = vieti - 1 & vieti = vieti - 1; \\%
}

\expandafter\newcommand\csname cqSyntaxRows@daca\endcsname{%
  daca (chei > 0) atunci & if (chei > 0) \{ & if chei > 0: & if (chei > 0) \{ \\
  \quad afiseaza "intra" & \quad cout << "intra"; & \quad print("intra") & \quad System.out.println("intra"); \\
  sfarsit & \} & & \} \\%
}

\expandafter\newcommand\csname cqSyntaxRows@cat-timp\endcsname{%
  cat timp (i <= 5) executa & while (i <= 5) \{ & while i <= 5: & while (i <= 5) \{ \\
  \quad i <- i + 1 & \quad i = i + 1; & \quad i = i + 1 & \quad i = i + 1; \\
  sfarsit & \} & & \} \\%
}

\expandafter\newcommand\csname cqSyntaxRows@pentru\endcsname{%
  pentru i de la 1 la 5 executa & for (int i = 1; i <= 5; i++) \{ & for i in range(1, 6): & for (int i = 1; i <= 5; i++) \{ \\
  \quad afiseaza i & \quad cout << i; & \quad print(i) & \quad System.out.println(i); \\
  sfarsit & \} & & \} \\%
}

\expandafter\newcommand\csname cqSyntaxRows@imbricate\endcsname{%
  pentru linie de la 1 la n executa & for (int linie = 1; linie <= n; linie++) \{ & for linie in range(1, n + 1): & for (int linie = 1; linie <= n; linie++) \{ \\
  \quad pentru coloana de la 1 la n executa & \quad for (int coloana = 1; coloana <= n; coloana++) \{ & \quad for coloana in range(1, n + 1): & \quad for (int coloana = 1; coloana <= n; coloana++) \{ \\
  \quad sfarsit & \quad \} & & \quad \} \\
  sfarsit & \} & & \} \\%
}

\newcommand{\SyntaxCheck}[1]{%
  \begin{languagebox}[Comparatie rapida: Pseudocod | C++ | Python | Java]
  \small
  \begin{tabularx}{\linewidth}{>{\ttfamily}X >{\ttfamily}X >{\ttfamily}X >{\ttfamily}X}
  \toprule
  Pseudocod & C++ & Python & Java \\
  \midrule
  \csname cqSyntaxRows@#1\endcsname
  \bottomrule
  \end{tabularx}
  \end{languagebox}
}

\makeatother
